\documentclass[12pt, a4paper]{article}

% --- Kodowanie i język ---
\usepackage[T1]{fontenc}
\usepackage[utf8]{inputenc}
\usepackage[polish]{babel}

% --- Czcionki ---
\usepackage{lmodern}
\usepackage{microtype}

% --- Marginesy ---
\usepackage[top=2.5cm, bottom=2.5cm, left=3cm, right=2.5cm]{geometry}

% --- Nagłówki i stopki ---
\usepackage{fancyhdr}
\pagestyle{fancy}
\fancyhf{}
\fancyhead[L]{\small Prywatność i Ochrona Danych Osobowych}
\fancyhead[R]{\small AI i Systemy Rekomendacyjne}
\fancyfoot[C]{\thepage}
\renewcommand{\headrulewidth}{0.4pt}

% --- Tabele ---
\usepackage{booktabs}
\usepackage{tabularx}
\usepackage{array}
\usepackage{longtable}
\usepackage{multirow}

% --- Kolory i ramki ---
\usepackage[table]{xcolor}
\definecolor{rowgray}{gray}{0.93}
\definecolor{headerblue}{RGB}{30, 70, 120}
\definecolor{accentblue}{RGB}{60, 120, 200}
\definecolor{lightblue}{RGB}{220, 235, 250}

% --- Hiperłącza ---
\usepackage[hidelinks, unicode]{hyperref}
\hypersetup{
  colorlinks=true,
  linkcolor=headerblue,
  urlcolor=accentblue,
  citecolor=headerblue
}

% --- Grafika i layout ---
\usepackage{graphicx}
\usepackage{float}
\usepackage{enumitem}
\usepackage{titlesec}
\usepackage{abstract}
\usepackage{caption}

% --- Formatowanie sekcji ---
\titleformat{\section}
  {\large\bfseries\color{headerblue}}
  {\thesection.}{0.6em}{}[\vspace{-4pt}\rule{\linewidth}{0.4pt}\vspace{2pt}]

\titleformat{\subsection}
  {\normalsize\bfseries\color{headerblue!80!black}}
  {\thesubsection.}{0.5em}{}

\titleformat{\subsubsection}
  {\normalsize\itshape\bfseries}
  {\thesubsubsection.}{0.5em}{}

% --- Odstępy ---
\setlength{\parindent}{1.2em}
\setlength{\parskip}{0.3em}
\usepackage{setspace}
\onehalfspacing

% --- Cytowania (thebibliography) ---

% --- Ramka na stronę tytułową ---
\usepackage{mdframed}

% =========================================================
\begin{document}

% =========================================================
%  STRONA TYTUŁOWA
% =========================================================
\begin{titlepage}
  \centering
  \vspace*{1.5cm}

  {\large\bfseries Projekt semestralny}\\[0.4cm]
  {\small Przedmiot: Prywatność i Ochrona Danych Osobowych}\\[2cm]

  \begin{mdframed}[linewidth=1.5pt, linecolor=headerblue,
    topline=true, bottomline=true, leftline=true, rightline=true,
    innertopmargin=16pt, innerbottommargin=16pt,
    innerrightmargin=20pt, innerleftmargin=20pt]
    \centering
    {\LARGE\bfseries\color{headerblue}
    Sztuczna inteligencja i systemy rekomendacyjne\\[0.4cm]
    a ochrona danych osobowych\\[0.3cm]}
    {\large Analiza przypadków Netflix, Instagram i Amazon\\
    w świetle wymogów RODO i zasady privacy-by-design}
  \end{mdframed}

  \vspace{2cm}

  {\normalsize Raport końcowy projektu}\\[0.3cm]

  \vfill
  {\small \today}
\end{titlepage}

% =========================================================
%  SPIS TREŚCI
% =========================================================
\tableofcontents
\newpage

% =========================================================
%  STRESZCZENIE
% =========================================================
\begin{abstract}
\noindent
Niniejszy raport stanowi końcowe podsumowanie projektu semestralnego poświęconego zagadnieniu sztucznej inteligencji i systemów rekomendacyjnych w kontekście ochrony danych osobowych. W pierwszej części omówiono kluczowe pojęcia teoretyczne: definicję i architekturę systemów rekomendacyjnych opartych na AI (AIRS), taksonomię zagrożeń prywatności oraz ramy prawne wynikające z RODO i zasady privacy-by-design. Następnie przeprowadzono szczegółową analizę trzech studiów przypadku: platformy streamingowej Netflix, platformy społecznościowej Instagram oraz platformy e-commerce Amazon. W kolejnej części dokonano komparatywnego zestawienia wszystkich trzech platform w zakresie minimalizacji danych, ograniczenia celu, przejrzystości, domyślnej ochrony prywatności i bezpieczeństwa technicznego. Raport kończy się wnioskami wskazującymi na systemowe napięcie między logiką maksymalizacji personalizacji a wymogami prawa ochrony danych.
\end{abstract}
\vspace{1em}
\hrule
\newpage

% =========================================================
%  CZĘŚĆ I – ZAGADNIENIA TEORETYCZNE
% =========================================================
\section{Zagadnienia teoretyczne: AI, systemy rekomendacyjne i prywatność}

\subsection{Systemy rekomendacyjne oparte na sztucznej inteligencji (AIRS)}

Systemy rekomendacyjne oparte na sztucznej inteligencji (ang.\ \textit{AI-based Recommender Systems}, AIRS) stanowią kluczowy element infrastruktury większości współczesnych platform cyfrowych. Ich podstawowym zadaniem jest przewidywanie, które produkty, treści, osoby lub usługi mogą zainteresować konkretnego użytkownika, oraz prezentowanie ich w sposób zwiększający prawdopodobieństwo interakcji lub zakupu \cite{chen2025, himeur2022}.

Z technicznego punktu widzenia AIRS mogą być klasyfikowane według kilku głównych paradygmatów \cite{himeur2022}. \textbf{Rekomendacje oparte na podobieństwie użytkowników} (\textit{collaborative filtering}) polegają na identyfikowaniu zbliżonych wzorców zachowania wśród różnych osób -- jeśli użytkownicy A i~B oglądali podobne treści, system zakłada, że treści lubiane przez A mogą spodobać się B. \textbf{Rekomendacje oparte na zawartości} (\textit{content-based filtering}) analizują cechy samych rekomendowanych elementów (gatunek, styl, atrybuty produktu) i dopasowują je do udokumentowanych preferencji użytkownika. \textbf{Modele hybrydowe} łączą oba podejścia, a najnowocześniejsze implementacje opierają się na głębokim uczeniu maszynowym i reprezentacjach wektorowych (\textit{embeddings}), pozwalając na uwzględnienie kontekstu, sekwencyjności zachowań i danych multimodalnych \cite{himeur2022, netflix_research}.

Efektywność AIRS jest ściśle uzależniona od ilości i jakości danych o użytkownikach \cite{chen2025}. Im więcej sygnałów behawioralnych platforma posiada, tym precyzyjniejsze mogą być jej predykcje. Ta fundamentalna zależność tworzy strukturalne napięcie z zasadami ochrony danych osobowych, w szczególności z zasadą minimalizacji danych i ograniczenia celu \cite{rodo, chen2025}.

\subsection{Taksonomia zagrożeń prywatności w systemach AIRS}

W literaturze przedmiotu wyróżnia się kilka kategorii zagrożeń prywatności specyficznych dla systemów rekomendacyjnych. Na podstawie badań jakościowych i ilościowych z udziałem ponad 580 użytkowników platform e-commerce zidentyfikowano trzy kluczowe grupy zagrożeń \cite{chen2025}:

\begin{description}[leftmargin=1.5em, style=nextline]
  \item[Poczucie bycia obserwowanym (\textit{perceived surveillance})]
    Użytkownicy odbierają rekomendacje jako dowód stałego nadzoru nad ich zachowaniem. Wzbudza to dyskomfort, szczególnie gdy rekomendacje wydają się wyprzedzać deklarowane potrzeby -- system „wie", czego użytkownik może zapragnąć, zanim sam to uświadomi.
  \item[Zagrożenie kradzieżą tożsamości (\textit{perceived identity theft})]
    Systemy AIRS przetwarzają i łączą dane identyfikacyjne z danymi behawioralnymi. W przypadku naruszenia bezpieczeństwa lub nieautoryzowanego dostępu taki profil może umożliwić nieuprawnionym podmiotom podszywanie się pod ofiarę lub popełnianie przestępstw na jej szkodę.
  \item[Nieuprawnione wtórne wykorzystanie danych (\textit{perceived unauthorized secondary use})]
    Dane zebrane w jednym celu (np.\ obsługa transakcji zakupowej, personalizacja treści) mogą być następnie wykorzystywane w celach, na które użytkownik nie wyraził świadomej i odrębnej zgody -- reklamie behawioralnej, profilowaniu dla podmiotów trzecich, trenowaniu modeli AI.
\end{description}

W literaturze wskazuje się również na zagrożenia od strony technicznej \cite{himeur2022}: ataki wstrzykiwania fałszywych profili (\textit{shilling attacks}), zatrucie danych treningowych (\textit{data poisoning}) oraz ryzyko reidentyfikacji pozornie zanonimizowanych zbiorów danych behawioralnych \cite{himeur2022, narayanan2008}. Ten ostatni problem ilustruje klasyczny przykład konkursu Netflix Prize z 2006~r., gdzie dowiedziono możliwości deanonimizacji zanonimizowanej bazy ocen filmów przy użyciu publicznych danych z serwisu IMDb \cite{narayanan2008}.

\subsection{Ramy prawne: RODO i zasada privacy-by-design}

Ogólne Rozporządzenie o Ochronie Danych (RODO, Rozporządzenie UE 2016/679) \cite{rodo} ustanawia zasady, które mają bezpośrednie znaczenie dla oceny systemów AIRS:

\begin{itemize}
  \item \textbf{Zasada legalności, rzetelności i przejrzystości} (art.~5 ust.~1 lit.~a RODO \cite{rodo}) -- dane muszą być przetwarzane w sposób rzetelny i transparentny dla osoby, której dotyczą. W kontekście systemów rekomendacyjnych oznacza to obowiązek zrozumiałego informowania użytkownika o tym, jak jego dane wpływają na prezentowane mu treści i oferty.
  \item \textbf{Zasada ograniczenia celu} (art.~5 ust.~1 lit.~b RODO \cite{rodo}) -- dane muszą być zbierane w konkretnych, wyraźnych i prawnie uzasadnionych celach oraz nie mogą być przetwarzane w sposób niezgodny z tymi celami. Personalizacja reklam prowadzona na podstawie danych zebranych pierwotnie w celu świadczenia usługi streamingowej lub realizacji transakcji zakupowej rodzi zasadnicze wątpliwości co do zgodności z tą zasadą.
  \item \textbf{Zasada minimalizacji danych} (art.~5 ust.~1 lit.~c RODO \cite{rodo}) -- przetwarzane dane muszą być adekwatne, stosowne i ograniczone do minimum niezbędnego do realizacji celów. Zasada ta pozostaje w bezpośrednim napięciu z logiką AIRS, które działają tym precyzyjniej, im więcej danych otrzymują \cite{chen2025}.
  \item \textbf{Zasada ograniczenia przechowywania} (art.~5 ust.~1 lit.~e RODO \cite{rodo}) -- dane nie powinny być przechowywane dłużej, niż jest to konieczne do realizacji celów.
  \item \textbf{Zasada integralności i poufności} (art.~5 ust.~1 lit.~f RODO \cite{rodo}) -- dane muszą być chronione przed nieuprawnionym dostępem i utratą.
\end{itemize}

Artykuł 25 RODO \cite{rodo} wprowadza koncepcję \textbf{ochrony danych w fazie projektowania i ochrony danych jako ustawienia domyślnego} (\textit{privacy-by-design} i \textit{privacy-by-default}). Zgodnie z wytycznymi Europejskiej Rady Ochrony Danych (EDPB) z 2019~r.\ (Guidelines 4/2019 on Article 25) \cite{edpb2019} zasada ta wymaga, aby ochrona prywatności była uwzględniana już na etapie projektowania systemu, nie dopiero po jego wdrożeniu. Privacy-by-default oznacza ponadto, że domyślne ustawienia usługi powinny ograniczać zakres przetwarzania do niezbędnego minimum -- użytkownik nie powinien musieć aktywnie wyłączać nadmiarowego profilowania \cite{edpb2019}.

Artykuły 15--22 RODO \cite{rodo} gwarantują podmiotom danych prawa dostępu, sprostowania, usunięcia, ograniczenia przetwarzania, przenoszenia, sprzeciwu oraz prawo do niepodlegania decyzji opartej wyłącznie na zautomatyzowanym przetwarzaniu (art.~22 RODO \cite{rodo}). To ostatnie prawo ma szczególne znaczenie w kontekście systemów AIRS, które automatycznie profilują użytkowników i na podstawie tego profilowania decydują o treściach, ofertach lub reklamach im prezentowanych.

Uzupełnieniem ram prawnych jest Akt o Usługach Cyfrowych (DSA, Rozporządzenie UE 2022/2065) \cite{dsa}, który wprowadza dla bardzo dużych platform internetowych szczegółowe obowiązki w zakresie przejrzystości systemów rekomendacyjnych. Art.~27 DSA \cite{dsa} zobowiązuje platformy do informowania użytkowników o głównych parametrach systemów rankingowych, a art.~38 DSA \cite{dsa} wymaga od bardzo dużych platform oferowania co najmniej jednej opcji rekomendacji nieopartej na profilowaniu.

\subsection{Pojęcie profilowania i jego znaczenie dla ochrony prywatności}

Profilowanie, zdefiniowane w art.~4 pkt~4 RODO \cite{rodo} jako forma zautomatyzowanego przetwarzania polegająca na ocenie czynników osobowych (zachowań, zainteresowań, lokalizacji, zdrowia, sytuacji ekonomicznej), jest centralnym mechanizmem działania AIRS. Z perspektywy ochrony danych profilowanie rodzi specyficzne ryzyka \cite{himeur2022, chen2025}:

\begin{itemize}
  \item \textbf{Inferencja danych wrażliwych} -- nawet jeśli użytkownik nie ujawnia wprost informacji o stanie zdrowia, orientacji seksualnej, poglądach politycznych czy przekonaniach religijnych, jego wzorce zachowania mogą umożliwiać statystyczne wnioskowanie o tych cechach \cite{himeur2022}. Historia oglądania, zakupów lub interakcji na platformie społecznościowej może tworzyć profil ujawniający dane szczególnej kategorii w rozumieniu art.~9 RODO \cite{rodo}.
  \item \textbf{Efekt bańki informacyjnej} -- algorytmiczne wzmacnianie treści zgodnych z dotychczasowymi zachowaniami może ograniczać ekspozycję użytkownika na odmienne perspektywy, treści neutralne lub sprzeczne z jego dotychczasowym profilem \cite{ig_ranking}.
  \item \textbf{Nieprzejrzystość decyzyjna} -- użytkownik rzadko dysponuje pełną wiedzą o tym, które konkretne dane i w jaki sposób wpłynęły na daną rekomendację, co ogranicza jego zdolność do korzystania z praw podmiotów danych \cite{chen2025, edpb2019}.
\end{itemize}

% =========================================================
%  CZĘŚĆ II – STUDIA PRZYPADKU
% =========================================================
\newpage
\section{Studia przypadku: analiza polityk prywatności}

\subsection{Netflix -- platforma streamingowa}

\subsubsection{Charakterystyka platformy i systemu rekomendacyjnego}

Netflix jest globalną platformą streamingową działającą w modelu subskrypcyjnym, umożliwiającą dostęp do filmów, seriali, programów dokumentalnych, animacji oraz gier. Personalizacja treści stanowi jeden z fundamentalnych elementów działania platformy -- użytkownik nie korzysta z jednego, neutralnego katalogu, lecz z interfejsu dynamicznie dopasowywanego do jego profilu i historii korzystania z usługi \cite{netflix_recs}.

System rekomendacyjny Netflixa szacuje prawdopodobieństwo, że dany tytuł będzie dla użytkownika interesujący, biorąc pod uwagę historię oglądania, oceny tytułów, podobieństwo do innych użytkowników, cechy treści, porę korzystania z usługi, preferowany język, typ urządzenia oraz czas oglądania \cite{netflix_recs, netflix_research}. Co istotne, personalizacja obejmuje wiele warstw interfejsu: wybór wierszy na stronie głównej, dobór tytułów w wierszu, kolejność pozycji oraz kontekst, w jakim są prezentowane \cite{netflix_recs}. System wpływa zatem nie tylko na to, \textit{co} jest proponowane, ale również \textit{w jakiej kolejności} i \textit{w jakim otoczeniu} to się pojawia.

\subsubsection{Zakres zbieranych danych osobowych}

Netflix przetwarza szerokie spektrum danych osobowych \cite{netflix_privacy, netflix_data}, które można podzielić na kilka kategorii:

\begin{enumerate}
  \item \textbf{Dane przekazywane przez użytkownika} -- adres e-mail, numer telefonu, dane logowania, informacje o płatności, dane konfiguracyjne profilu (nazwa, ikona, preferencje językowe, lista ``Moja lista'', ustawienia napisów, blokady wiekowe, oceny treści) \cite{netflix_privacy}.
  \item \textbf{Dane zbierane automatycznie} -- zdarzenia odtwarzania (rozpoczęcie, zatrzymanie, pauza, zakończenie), historia oglądania, historia gier, wyszukiwane frazy, kliknięcia, wyświetlenia stron, czas dostępu, reakcje na tytuły, długość sesji \cite{netflix_data, netflix_recs}.
  \item \textbf{Dane urządzeń i sieci} -- identyfikatory urządzeń, adresy IP, typ urządzenia, system operacyjny, dane przeglądarki, typ połączenia, pliki cookie, identyfikatory reklamowe \cite{netflix_privacy}.
  \item \textbf{Dane od partnerów} -- dane pozyskiwane od producentów telewizorów, dostawców Internetu, operatorów telefonii, partnerów rozliczeniowych i dostawców urządzeń streamingowych \cite{netflix_privacy}.
  \item \textbf{Dane reklamowe} -- informacje o reklamach wyświetlanych użytkownikowi, interakcjach z reklamami, identyfikatorach urządzeń oraz dane od firm reklamowych (dotyczy planów subskrypcyjnych z reklamami) \cite{netflix_privacy}.
\end{enumerate}

\subsubsection{Identyfikacja zagrożeń prywatności}

Analiza przypadku Netflixa pozwala wyodrębnić kilka kluczowych zagrożeń dla prywatności użytkowników \cite{netflix_privacy, netflix_data}. \textbf{Stałe profilowanie behawioralne} -- analiza nie tylko jawnych ocen, ale też czasu oglądania, momentów pauzy, tempa przewijania i pory korzystania -- może tworzyć profil ujawniający pośrednio zainteresowania zdrowotne, polityczne, religijne lub emocjonalne użytkownika \cite{chen2025}. Problem \textbf{nadmiernej personalizacji} polega na tym, że algorytm wzmacniający określone kategorie treści może prowadzić do zawężenia faktycznego wyboru, mimo formalnego dostępu do całego katalogu \cite{netflix_recs, netflix_research}.

Specyficznym ryzykiem jest \textbf{prywatność wewnątrz konta} -- profile zapewniają personalizację, ale nie zawsze pełną separację prywatności między osobami korzystającymi z jednego konta \cite{netflix_transfer}. Historyczny przypadek \textbf{Netflix Prize} \cite{narayanan2008} pozostaje modelowym przykładem ryzyka reidentyfikacji: zanonimizowana baza danych o preferencjach filmowych okazała się możliwa do powiązania z konkretnymi osobami przy pomocy zewnętrznych źródeł. Wreszcie, \textbf{wtórne wykorzystanie danych} do reklam behawioralnych i marketingu na usługach stron trzecich rodzi pytania o granice pierwotnego celu przetwarzania \cite{netflix_privacy, rodo}.

\subsubsection{Ocena zgodności z RODO i zasadą privacy-by-design}

\textbf{Minimalizacja danych} -- Netflix pozytywnie wyróżnia się deklaracją, że system rekomendacyjny nie wykorzystuje bezpośrednio danych demograficznych (wieku, płci) w procesie decyzyjnym \cite{netflix_recs}. Jednocześnie zakres danych behawioralnych jest bardzo szeroki i wykracza poza to, co jest bezwarunkowo konieczne do rekomendacji treści, szczególnie w kontekście danych reklamowych i partnerskich \cite{netflix_privacy, rodo}.

\textbf{Ograniczenie celu} -- podstawowym celem jest świadczenie usługi streamingowej, ale te same lub podobne dane są wykorzystywane do marketingu, reklam i integracji partnerskich \cite{netflix_privacy}. Granice między tymi celami mogą być dla użytkownika niejasne \cite{chen2025}.

\textbf{Przejrzystość} -- Netflix wypada stosunkowo dobrze: platforma posiada osobny artykuł wyjaśniający działanie systemu rekomendacyjnego prostym językiem \cite{netflix_recs}. Nadal jednak użytkownik nie widzi konkretnego wyjaśnienia przy każdej rekomendacji ani pełnego profilu przypisanych mu cech \cite{edpb2019}.

\textbf{Privacy-by-default} -- personalizacja jest domyślna i stanowi rdzeń usługi \cite{netflix_recs}. Istnieją kontrole (profile dziecięce, blokada PIN, ukrywanie tytułów z historii) \cite{netflix_stop}, ale przeciętny użytkownik musi aktywnie ich szukać, co stoi w sprzeczności z wymogami art.~25 ust.~2 RODO \cite{rodo, edpb2019}.

\textbf{Prawa podmiotów danych} -- Netflix zapewnia mechanizmy dostępu, aktualizacji, usunięcia danych i pobrania kopii \cite{netflix_data}. Użytkownik może korzystać z ustawień prywatności, zarządzać historią oglądania i kontrolować reklamy behawioralne \cite{netflix_stop}.

\subsubsection{Rekomendacje}

Na podstawie analizy można wskazać następujące kierunki poprawy zgodności z privacy-by-design: wyraźne rozdzielenie personalizacji treści od personalizacji reklam, panel kontroli profilu rekomendacyjnego z możliwością usuwania poszczególnych sygnałów, łatwiejsze resetowanie historii rekomendacyjnej, silniejsza domyślna separacja profili na kontach współdzielonych, skrócenie retencji surowych danych behawioralnych oraz stosowanie pseudonimizacji i agregacji w celach analitycznych.

% ---------------------------------------------------------
\subsection{Instagram -- platforma społecznościowa}
% ---------------------------------------------------------

\subsubsection{Charakterystyka platformy i systemu rekomendacyjnego}

Instagram jest platformą społecznościową należącą do Meta Platforms, służącą do publikowania, oglądania i udostępniania zdjęć, filmów, relacji, rolek i wiadomości prywatnych. W odróżnieniu od Netflixa, gdzie rekomendacja dotyczy głównie wyboru treści rozrywkowych, system rekomendacyjny Instagrama wpływa na obieg informacji społecznej: decyduje o widoczności postów znajomych, popularności twórców, ekspozycji treści politycznych i skuteczności reklam \cite{ig_ranking}.

Platforma nie stosuje jednego algorytmu -- różne części aplikacji (\textbf{Feed}, \textbf{Stories}, \textbf{Explore}, \textbf{Reels}) posiadają własne systemy rankingowe \cite{ig_ranking, ig_creators}. System analizuje wcześniejszą aktywność użytkownika, polubienia, komentarze, zapisane treści, udostępnienia, wyszukiwania, relacje z innymi kontami, historię interakcji z autorami oraz informacje o samych postach \cite{ig_ranking, ig_creators}.

\subsubsection{Zakres zbieranych danych osobowych}

Meta przetwarza w ramach Instagrama następujące kategorie danych \cite{meta_privacy}:

\begin{enumerate}
  \item \textbf{Dane przekazywane przez użytkownika} -- imię i nazwisko, nazwa profilu, adres e-mail, numer telefonu, zdjęcie profilowe, data urodzenia, treści publikowanych postów, relacji i rolek \cite{meta_privacy}.
  \item \textbf{Dane o aktywności} -- polubienia, komentarze, zapisane posty, udostępnienia, obejrzane rolki, czas oglądania, wyszukiwania, kliknięcia w profile, reakcje na Stories, tempo przewijania, ukrywanie treści, obserwowanie kont \cite{meta_privacy, ig_creators}.
  \item \textbf{Dane o relacjach społecznych} -- lista obserwowanych i obserwujących, częstotliwość kontaktu z konkretnymi osobami, oznaczenia, wiadomości \cite{meta_privacy}.
  \item \textbf{Dane techniczne i lokalizacyjne} -- urządzenie, system operacyjny, adres IP, pliki cookie, identyfikatory urządzeń, język, strefa czasowa, parametry sieci \cite{meta_privacy}.
  \item \textbf{Dane reklamowe i dane spoza platformy} -- dane od reklamodawców, partnerów biznesowych, właścicieli stron korzystających z narzędzi Meta (piksela Meta, SDK) \cite{meta_privacy}.
  \item \textbf{Dane z funkcji AI} -- interakcje użytkowników z funkcjami generatywnej AI mogą być wykorzystywane do personalizacji treści i reklam na platformach Meta \cite{meta_ai}.
\end{enumerate}

\subsubsection{Identyfikacja zagrożeń prywatności}

Instagram generuje szczególnie rozbudowany katalog zagrożeń prywatności. \textbf{Inferencja danych wrażliwych} -- regularne przeglądanie treści dotyczących zdrowia, polityki, religii lub seksualności może prowadzić do wnioskowania o danych szczególnie chronionych w rozumieniu art.~9 RODO \cite{rodo}, nawet jeśli użytkownik nigdy ich wprost nie ujawnił \cite{chen2025, himeur2022}. \textbf{Nadmierna personalizacja i bańka informacyjna} -- stopniowe wzmacnianie treści podobnych do dotychczasowych interakcji może ograniczać ekspozycję na odmienne poglądy i informacje \cite{ig_ranking}; w kontekście platformy społecznościowej skutki są poważniejsze niż przy wyborze filmów. \textbf{Manipulacja uwagą i dark patterns} -- system optymalizowany pod kątem zaangażowania rodzi pytania o autonomię użytkownika; w 2025~r.\ holenderski sąd nakazał Meta zmianę domyślnych ustawień osi czasu Facebooka i Instagrama \cite{reuters_nl}. \textbf{Wtórne wykorzystanie danych} -- dane o aktywności wpływają jednocześnie na rekomendacje treści, personalizację reklam, analizę skuteczności kampanii i ulepszanie modeli AI \cite{meta_privacy, meta_ai}. \textbf{Wpływ na małoletnich} -- nastolatkowie są szczególnie podatni na treści dotyczące wyglądu, popularności i zdrowia psychicznego \cite{meta_teen}.

W 2026~r.\ irlandzki regulator medialny wszczął postępowania dotyczące manipulacji algorytmicznymi feedami \cite{reuters_ie}, a Komisja Europejska wstępnie stwierdziła, że Meta narusza DSA przez brak skutecznych mechanizmów zapobiegania dostępowi dzieci poniżej 13. roku życia do Instagrama i Facebooka \cite{ec_meta_dsa}.

\subsubsection{Ocena zgodności z RODO i zasadą privacy-by-design}

\textbf{Minimalizacja danych} -- zakres przetwarzanych danych jest bardzo szeroki, szczególnie w obszarze danych behawioralnych i reklamowych \cite{meta_privacy}. Dane z aktywności w Instagramie mogą łączyć się z danymi z innych produktów Meta, co znacznie rozbudowuje profil użytkownika \cite{meta_privacy, himeur2022}.

\textbf{Ograniczenie celu} -- cele przetwarzania obejmują świadczenie usługi, personalizację treści, reklamy, bezpieczeństwo, analitykę, rozwój produktów i integrację z innymi usługami Meta \cite{meta_privacy}. Te same dane mogą jednocześnie służyć wielu celom, co utrudnia weryfikację zgodności z zasadą ograniczenia celu \cite{rodo, chen2025}.

\textbf{Przejrzystość} -- Instagram publikuje ogólne wyjaśnienia działania rankingów \cite{ig_ranking, ig_creators}. Brak jednak szczegółowych, kontekstowych wyjaśnień przy konkretnych rekomendacjach \cite{edpb2019}. Ustawienia prywatności są rozproszone w różnych panelach \cite{ig_recs}.

\textbf{Privacy-by-default} -- personalizacja jest domyślna \cite{ig_ranking}. Dostępne są mechanizmy kontroli (reset sugerowanych treści \cite{ig_reset}, ustawienia wrażliwych treści, konta nastoletnie \cite{meta_teen}), ale ich znalezienie wymaga aktywnych działań użytkownika, co jest sprzeczne z wymogami art.~25 RODO \cite{rodo, edpb2019}.

\textbf{DSA} -- jako bardzo duża platforma, Instagram podlega szczegółowym wymogom DSA dotyczącym przejrzystości systemów rekomendacyjnych i oferowania opcji feedu nieopartego na profilowaniu \cite{dsa, ec_dsa}.

\subsubsection{Rekomendacje}

Analiza sugeruje następujące kierunki poprawy: wyraźne rozdzielenie rekomendacji treści od personalizacji reklam, panel profilu rekomendacyjnego z możliwością usuwania przypisanych kategorii zainteresowań, łatwy i trwały wybór feedu nieopartego na profilowaniu, kontekstowe wyjaśnienia ``dlaczego to widzę?'', ograniczenie retencji szczegółowych danych behawioralnych, silniejsza ochrona przed inferencją danych wrażliwych, domyślnie bezpieczniejsze ustawienia dla małoletnich oraz niezależne audyty systemów rekomendacyjnych.

% ---------------------------------------------------------
\subsection{Amazon -- platforma e-commerce}
% ---------------------------------------------------------

\subsubsection{Charakterystyka platformy i systemu rekomendacyjnego}

Amazon jest jedną z największych platform e-commerce na świecie i pionierem stosowania systemów rekomendacyjnych w handlu elektronicznym. Według powszechnie cytowanych szacunków rekomendacje odpowiadają za około 35\% zakupów dokonywanych na platformie. Stosowane przez Amazon rozwiązanie oparte na filtrowaniu kolaboratywnym zorientowanym na przedmioty (\textit{item-to-item collaborative filtering}) stało się wzorcem dla całej branży \cite{himeur2022}.

Amazon wdraża mechanizmy rekomendacyjne w wielu punktach styku z użytkownikiem: na stronie głównej, na kartach produktów, w procesie zakupu, w wiadomościach e-mail oraz w sekcji ``Klienci, którzy kupili ten produkt, kupili również''. Do obsługi silników personalizacji służy m.in.\ Amazon Personalize -- własna platforma dostępna w ramach AWS \cite{amazon_dsp}. System zasilany jest danymi ze wszystkich obszarów ekosystemu Amazona: platformy zakupowej, urządzeń Kindle, serwisu Prime Video, asystenta głosowego Alexa oraz -- dla posiadaczy kart lojalnościowych -- sieci sklepów Whole Foods \cite{amazon_privacy}.

\subsubsection{Zakres zbieranych danych osobowych}

Amazon gromadzi dane w trzech zasadniczych kategoriach \cite{amazon_privacy}:

\begin{enumerate}
  \item \textbf{Dane przekazywane przez użytkownika} -- imię i nazwisko, adres, numery telefonów, informacje płatnicze, wiek, dane logowania, zdjęcia profilowe, nagrania głosowe (asystent Alexa i urządzenia Echo) \cite{amazon_privacy}.
  \item \textbf{Dane zbierane automatycznie} -- adres IP, historia wyszukiwań, historia przeglądania kart produktów, historia zakupów, czas spędzony na stronach, liczba i czas sesji streamingowych, dane o interakcjach (kliknięcia, przewijanie, najechania kursorem), dane urządzenia, lokalizacja geograficzna \cite{amazon_privacy}.
  \item \textbf{Dane z zewnętrznych źródeł} -- zaktualizowane informacje adresowe od firm kurierskich, dane o interakcjach z partnerami handlowymi, informacje z biur informacji kredytowej \cite{amazon_privacy}.
\end{enumerate}

Amazon wprost deklaruje, że dane osobowe są wykorzystywane do rekomendowania funkcji, produktów i usług potencjalnie interesujących użytkownika oraz do personalizowania doświadczenia \cite{amazon_privacy}. Zakres gromadzonych informacji czyni Amazon jednym z podmiotów o najintensywniejszym profilowaniu użytkowników na rynku cyfrowym \cite{chen2025, himeur2022}.

\subsubsection{Identyfikacja zagrożeń prywatności}

Analiza Amazona w kontekście opisanej taksonomii zagrożeń \cite{chen2025} wskazuje trzy główne kategorie. \textbf{Poczucie bycia obserwowanym} -- Amazon zbiera dane jednocześnie z wielu niezależnych źródeł (platforma zakupowa, reklamy na stronach partnerskich korzystających z Amazon Advertising, historia głosowa z Alexy, zachowania w sklepach Whole Foods), co może prowadzić do rekomendacji sprawiających wrażenie, że platforma ``zna'' potrzeby użytkownika zanim on sam je uświadomi.

\textbf{Zagrożenie kradzieżą tożsamości} -- Amazon przetwarza dane wrażliwe: numery kart kredytowych, adresy zamieszkania, dane rozliczeń finansowych, a w przypadku sprzedawców -- numery identyfikacji podatkowej i informacje bankowe. Połączenie danych identyfikacyjnych z danymi behawioralnymi tworzy pełny profil osoby, który w razie naruszenia bezpieczeństwa może być użyty do popełnienia przestępstw.

\textbf{Nieuprawnione wtórne wykorzystanie danych} -- dane zakupowe mogą być przekazywane firmom analitycznym, agencjom reklamowym i partnerom płatności. Platforma Amazon DSP (\textit{Demand-Side Platform}) \cite{amazon_dsp} umożliwia wyświetlanie reklam opartych na historii zachowań na Amazonie na tysiącach innych witryn -- model \textit{cross-site behavioral advertising} powoduje, że dane zebrane w celach transakcyjnych są intensywnie wykorzystywane w odmiennych kontekstach \cite{amazon_dsp, rodo}.

\subsubsection{Ocena zgodności z RODO i zasadą privacy-by-design}

\textbf{Minimalizacja danych} -- zakres gromadzonych informacji, obejmujący dane głosowe, aktywność w sklepach stacjonarnych i historię przeglądania stron partnerskich, znacznie wykracza poza to, co byłoby niezbędne do realizacji transakcji zakupowej \cite{amazon_privacy, rodo}. Jest to najpoważniejsze uchybienie Amazona względem zasad RODO.

\textbf{Ograniczenie celu} -- praktyka udostępniania danych behawioralnych partnerom reklamowym przez Amazon DSP \cite{amazon_dsp} budzi istotne wątpliwości co do zgodności z art.~5 ust.~1 lit.~b RODO \cite{rodo}, szczególnie z perspektywy użytkownika dokonującego zakupu.

\textbf{Przejrzystość} -- polityka prywatności jest szczegółowa i dostępna publicznie \cite{amazon_privacy}, jednak jej objętość i poziom skomplikowania sprawiają, że przeciętny użytkownik nie jest w stanie zrozumieć pełnego zakresu operacji przetwarzania \cite{chen2025}.

\textbf{Privacy-by-default} -- użytkownik, który nie podejmie aktywnych działań (zarządzanie preferencjami reklam w sekcji ``Your Ads Privacy Choices'', konfiguracja ustawień konta), automatycznie podlega pełnemu zakresowi personalizacji i profilowania \cite{amazon_privacy}. Domyślna konfiguracja odpowiada maksymalizacji przetwarzania, nie jego minimalizacji, co stoi w sprzeczności z art.~25 ust.~2 RODO \cite{rodo, edpb2019}.

\textbf{Bezpieczeństwo techniczne} -- w tym obszarze Amazon prezentuje wysoki poziom dojrzałości: certyfikacja PCI DSS, szyfrowanie danych w chmurze AWS, wielostopniowe mechanizmy uwierzytelniania \cite{amazon_privacy}. Zaawansowanie techniczne w zakresie bezpieczeństwa nie jest jednak równoznaczne z wdrożeniem privacy-by-design \cite{edpb2019}.

\subsubsection{Kontekst regulacyjny}

W 2021~r.\ luksemburska Komisja Ochrony Danych nałożyła na Amazon karę za przetwarzanie danych do celów reklamy behawioralnej w sposób uznany za niezgodny z RODO \cite{rodo}. W marcu 2026~r.\ luksemburski sąd uchylił nałożoną karę i skierował sprawę do ponownej oceny, jednak samo postępowanie ilustruje, że praktyki profilowania Amazona budziły poważne wątpliwości organów nadzorczych \cite{amazon_privacy}.

% =========================================================
%  CZĘŚĆ III – PORÓWNANIE
% =========================================================
\newpage
\section{Porównanie polityk prywatności: Netflix, Instagram, Amazon}

\subsection{Cel i zakres porównania}

Celem niniejszej części jest zestawienie wszystkich trzech platform pod kątem kluczowych kryteriów ochrony prywatności: minimalizacji danych, ograniczenia celu, przejrzystości, privacy-by-default, bezpieczeństwa technicznego oraz ogólnego poziomu ryzyka. Platforma streamingowa, społecznościowa i e-commerce reprezentują trzy odmienne modele biznesowe, co wpływa na charakter i zakres przetwarzania danych, a tym samym na naturę i intensywność zagrożeń dla prywatności.

\subsection{Syntetyczne zestawienie porównawcze}

\begin{longtable}{>{\bfseries\small}p{3.4cm}
                  >{\small\raggedright}p{3.4cm}
                  >{\small\raggedright}p{3.4cm}
                  >{\small\raggedright\arraybackslash}p{3.4cm}}
  \toprule
  \rowcolor{headerblue}
  \textcolor{white}{Kryterium} &
  \textcolor{white}{Netflix} &
  \textcolor{white}{Instagram / Meta} &
  \textcolor{white}{Amazon} \\
  \midrule
  \endfirsthead
  \toprule
  \rowcolor{headerblue}
  \textcolor{white}{Kryterium} &
  \textcolor{white}{Netflix} &
  \textcolor{white}{Instagram / Meta} &
  \textcolor{white}{Amazon} \\
  \midrule
  \endhead
  \bottomrule
  \endfoot

  \rowcolor{lightblue}
  Typ platformy &
  Streamingowa (subskrypcja) &
  Społecznościowa (reklamy) &
  E-commerce i ekosystem usług \\

  Główne dane do rekomendacji &
  Historia oglądania, oceny, czas oglądania, urządzenie &
  Aktywność w aplikacji, relacje społeczne, dane z produktów Meta &
  Historia zakupów, przeglądania, wyszukiwań, dane z całego ekosystemu \\

  \rowcolor{lightblue}
  Minimalizacja danych &
  Umiarkowane ryzyko -- zakres uzasadniony streaming, problematyczny przy reklamach &
  Wysokie ryzyko -- bardzo szeroki zakres, łączenie danych między usługami Meta &
  Najwyższe ryzyko -- dane głosowe, sklepowe, partnerskie wykraczają poza cel zakupowy \\

  Ograniczenie celu &
  Ryzyko głównie przy reklamach behawioralnych i partnerach &
  Wysokie ryzyko -- cel rekomendacji miesza się z celami reklamowymi i analizą AI &
  Wysokie ryzyko -- dane zakupowe intensywnie wykorzystywane do reklamy cross-site (DSP) \\

  \rowcolor{lightblue}
  Przejrzystość &
  Lepsza -- osobny artykuł o zasadach działania rekomendacji &
  Ogólna -- wyjaśnione istnienie różnych systemów, brak kontekstowych wyjaśnień &
  Ograniczona -- polityka szczegółowa, ale ekosystem zbyt złożony dla przeciętnego użytkownika \\

  Privacy-by-default &
  Częściowo -- personalizacja domyślna, ale istnieją kontrole &
  Problematyczne -- personalizacja domyślna, opcje trudno dostępne &
  Najbardziej problematyczne -- domyślna konfiguracja maksymalizuje profilowanie \\

  \rowcolor{lightblue}
  Bezpieczeństwo techniczne &
  Dobre -- zarządzanie urządzeniami, PIN, szyfrowanie &
  Dobre -- weryfikacja logowań, mechanizmy zgłaszania naruszeń &
  Bardzo dobre -- PCI DSS, AWS, szyfrowanie end-to-end, uwierzytelnianie wieloskładnikowe \\

  Profilowanie behawioralne &
  Wysokie w zakresie gustów oglądania &
  Bardzo wysokie -- obejmuje sieć społeczną, emocje, zainteresowania &
  Bardzo wysokie -- obejmuje zachowania zakupowe, finansowe, głosowe, fizyczne \\

  \rowcolor{lightblue}
  Ryzyko inferencji danych wrażliwych &
  Średnie -- historia oglądania może ujawniać cechy szczególne &
  Wysokie -- aktywność w aplikacji może ujawniać zdrowie, poglądy, seksualność &
  Wysokie -- profil zakupowy może ujawniać stan zdrowia, sytuację finansową, przekonania \\

  Wpływ na użytkownika &
  Wybór treści, czas korzystania z platformy &
  Obieg informacji, zdrowie psychiczne, poglądy, relacje &
  Decyzje zakupowe, wydatki, ekspozycja na reklamy poza platformą \\

  \rowcolor{lightblue}
  Kontekst regulacyjny &
  RODO; postępowania dotyczące reklam &
  RODO, DSA; postępowania dot.\ manipulacji feedem i ochrony małoletnich &
  RODO; postępowanie CNPD i wyrok sądu w 2026~r.\ \\

  Ogólna ocena ryzyka &
  Średnie do wysokiego &
  Wysokie &
  Wysokie do bardzo wysokiego \\
\end{longtable}

\subsection{Analiza porównawcza według kryteriów RODO}

\subsubsection{Minimalizacja danych}

Spośród trzech platform najlepiej z zasadą minimalizacji radzi sobie Netflix -- zakres zbieranych danych jest szeroki, ale przynajmniej jego rdzeń (historia oglądania, oceny, profil) jest bezpośrednio powiązany z główną funkcją usługi \cite{netflix_data, netflix_recs}. Zakres danych Amazona stanowi najpoważniejsze wyzwanie: dane głosowe zbierane przez Alexę, zachowania w sklepach Whole Foods, historia przeglądania na stronach partnerskich -- wszystkie te elementy wykraczają daleko poza to, co byłoby niezbędne do realizacji transakcji zakupowej \cite{amazon_privacy, rodo}. Instagram zajmuje pozycję pośrednią, choć szczególnie problematyczna jest możliwość łączenia danych z wielu produktów Meta i wzbogacania profilu o aktywność użytkownika poza samą aplikacją \cite{meta_privacy, himeur2022}.

\subsubsection{Ograniczenie celu}

We wszystkich trzech przypadkach obserwujemy to samo systemowe zjawisko: dane zbierane pierwotnie do określonego celu (streaming, interakcja społecznościowa, zakup) są intensywnie wykorzystywane do celów reklamowych \cite{rodo, chen2025}. Amazon jest pod tym względem najbardziej problematyczny ze względu na Amazon DSP i skalę reklamy behawioralnej poza platformą \cite{amazon_dsp}. Instagram jest trudny do oceny ze względu na wielość równoległych celów: te same dane o aktywności mogą jednocześnie napędzać rekomendacje Reels, personalizację reklam, ulepszanie modeli AI i pomiar skuteczności kampanii reklamodawców \cite{meta_privacy, meta_ai}. Netflix, dzięki subskrypcyjnemu modelowi biznesowemu, ma mniejszy bezwzględny imperatyw do wtórnego wykorzystania danych, choć rozwój planów reklamowych zmienia tę sytuację \cite{netflix_privacy}.

\subsubsection{Przejrzystość}

Netflix wyróżnia się pozytywnie: platforma w sposób dostępny dla laika opisuje podstawowe zasady działania systemu rekomendacyjnego \cite{netflix_recs, netflix_research}. Instagram informuje o istnieniu różnych systemów rankingowych \cite{ig_ranking, ig_creators}, ale brakuje kontekstowych, indywidualnych wyjaśnień. Amazon jest najmniej przejrzysty w praktycznym wymiarze -- polityka prywatności jest formalnie kompletna \cite{amazon_privacy}, ale złożoność ekosystemu (AWS, Alexa, Prime Video, Whole Foods, partnerzy reklamowi, sprzedawcy zewnętrzni) sprawia, że przeciętny użytkownik nie jest w stanie zrozumieć, jak i w jakim celu jego dane są łączone. Przejrzystość realizowana wyłącznie na poziomie dokumentu prawnego nie przekłada się automatycznie na przejrzystość funkcjonalną z perspektywy użytkownika \cite{chen2025}.

\subsubsection{Privacy-by-default}

Żadna z analizowanych platform nie spełnia w pełni wymagań art.~25 ust.~2 RODO \cite{rodo, edpb2019}, choć poziom uchybień jest zróżnicowany. Netflix i Instagram traktują szeroko zakrojoną personalizację jako podstawową cechę usługi \cite{netflix_recs, ig_ranking}, co jest biznesowo racjonalne, ale oznacza, że użytkownik od pierwszego logowania podlega intensywnemu profilowaniu. Amazon jest najbardziej odległy od ideału privacy-by-default -- domyślna konfiguracja konta automatycznie włącza reklamę behawioralną, profilowanie wielokanałowe i udostępnianie danych partnerom \cite{amazon_privacy, amazon_dsp}. Szczególnie istotna jest różnica między dostępnością narzędzi kontroli a ich domyślnymi ustawieniami: wszystkie trzy platformy oferują mechanizmy ograniczania przetwarzania \cite{netflix_stop, ig_reset, amazon_privacy}, ale żadna nie uczyniła ochrony prywatności domyślnym stanem systemu \cite{edpb2019}.

\subsubsection{Bezpieczeństwo techniczne}

Pod względem bezpieczeństwa technicznego najbardziej dojrzały jest Amazon -- certyfikacja PCI DSS, szyfrowanie w chmurze AWS i wielostopniowe uwierzytelnianie odzwierciedlają poważne podejście do ochrony danych finansowych i transakcyjnych \cite{amazon_privacy}. Netflix i Instagram również stosują zaawansowane mechanizmy bezpieczeństwa, dostosowane do specyfiki swoich usług \cite{netflix_privacy, meta_privacy}. Należy jednak podkreślić, że bezpieczeństwo techniczne jest warunkiem koniecznym, ale nie wystarczającym dla zgodności z privacy-by-design \cite{edpb2019}: system może doskonale chronić dane przed nieautoryzowanym dostępem, a jednocześnie zbierać nadmierny ich zakres lub używać ich do celów sprzecznych z oczekiwaniami użytkowników \cite{himeur2022}.

\subsection{Przekrojowe wnioski z analizy porównawczej}

Porównanie trzech platform ujawnia kilka istotnych prawidłowości. Po pierwsze, \textbf{model biznesowy determinuje intensywność zagrożeń prywatności} \cite{chen2025}: platforma finansowana z reklam (Instagram) ma strukturalną zachętę do zbierania możliwie dużo danych o zainteresowaniach i zachowaniach użytkowników \cite{meta_privacy}; platforma e-commerce (Amazon) łączy dane transakcyjne, behawioralne i lokalizacyjne, tworząc profile o bezprecedensowej głębokości \cite{amazon_privacy}; platforma subskrypcyjna (Netflix) ma nieco mniejszą bezpośrednią zachętę do ekstensywnego profilowania, choć rozwój planów reklamowych zmienia tę dynamikę \cite{netflix_privacy}.

Po drugie, \textbf{skala ekosystemu potęguje ryzyko} \cite{himeur2022}: Amazon i Instagram działają jako węzły rozległych ekosystemów (odpowiednio AWS/Alexa/Whole Foods i Meta/Facebook/Messenger), co oznacza, że dane użytkownika są łączone z danymi z wielu niezależnych punktów styku \cite{amazon_privacy, meta_privacy}. Netflix jest pod tym względem relatywnie bardziej ograniczony \cite{netflix_privacy}.

Po trzecie, \textbf{wpływ rekomendacji na użytkownika rośnie wraz ze znaczeniem domeny}: rekomendacje filmów wpływają na czas wolny \cite{netflix_recs}; rekomendacje na platformie społecznościowej wpływają na obraz świata, relacje i zdrowie psychiczne \cite{ig_ranking, reuters_nl}; rekomendacje zakupowe wpływają bezpośrednio na decyzje finansowe \cite{amazon_privacy}. Ocena ryzyka prywatności powinna uwzględniać nie tylko zakres przetwarzania, ale też konsekwencje, jakie to przetwarzanie niesie dla życia użytkownika \cite{chen2025}.

% =========================================================
%  CZĘŚĆ IV – WNIOSKI
% =========================================================
\newpage
\section{Wnioski końcowe}

\subsection{Systemowe napięcie między personalizacją a ochroną prywatności}

Przeprowadzone analizy potwierdzają tezę, którą można sformułować jako fundamentalne napięcie: \textbf{ekonomiczna logika nowoczesnych systemów rekomendacyjnych AIRS jest strukturalnie sprzeczna z podstawowymi zasadami ochrony danych osobowych} \cite{chen2025, himeur2022}. Im więcej danych o użytkowniku system posiada, tym skuteczniej może przewidywać jego preferencje, wydłużać czas spędzany na platformie, zwiększać wskaźniki konwersji i podnosić wartość reklam. Tymczasem zasady minimalizacji danych, ograniczenia celu i domyślnej ochrony prywatności \cite{rodo, edpb2019} zmierzają w dokładnie przeciwnym kierunku -- ku ograniczeniu zakresu przetwarzania i jego celów.

To napięcie nie jest możliwe do całkowitego usunięcia: pewien poziom profilowania jest inherentny dla jakiegokolwiek systemu rekomendacyjnego. Problemem jest nie samo profilowanie, lecz jego zakres, przejrzystość, kontrolowalność przez użytkownika i powiązanie z wtórnymi celami reklamowymi i analitycznymi, które wykraczają poza pierwotną funkcję usługi.

\subsection{Privacy-by-design jako wymóg systemowy, nie deklaratywny}

Kluczowy wniosek z porównania wszystkich trzech platform brzmi: \textbf{privacy-by-design nie może być rozumiane jako zestaw zabezpieczeń technicznych dodawanych do gotowego systemu ani jako deklaracja zamieszczana w polityce prywatności} \cite{edpb2019}. Musi stanowić zasadę organizującą cały cykl życia systemu -- od wyboru architektury przetwarzania, przez dobór sygnałów używanych przez algorytm, po projektowanie interfejsu użytkownika i domyślnych ustawień \cite{edpb2019, rodo}.

W praktyce żadna z analizowanych platform nie spełnia w pełni tego ideału. Netflix jest w tej stawce relatywnie najlepszy -- jego polityka prywatności i dokumentacja systemu rekomendacyjnego są przejrzyste \cite{netflix_privacy, netflix_recs}, a zakres danych jest przynajmniej częściowo powiązany z podstawową funkcją usługi \cite{netflix_data}. Amazon i Instagram mają poważniejsze uchybienia: w obu przypadkach domyślne ustawienia maksymalizują profilowanie \cite{amazon_privacy, meta_privacy}, zakres danych wykracza daleko poza to, co jest niezbędne, a przepływ danych między różnymi usługami i partnerami jest trudny do pełnego zrozumienia przez użytkownika \cite{chen2025}.

\subsection{Prawa podmiotów danych w erze AIRS}

Analiza pokazuje, że formalna dostępność mechanizmów realizacji praw podmiotów danych (dostęp, sprostowanie, usunięcie, sprzeciw) nie jest równoznaczna z realną możliwością korzystania z tych praw w kontekście systemów rekomendacyjnych \cite{chen2025, edpb2019}. Użytkownik, który chce zrozumieć, dlaczego widzi konkretne rekomendacje, rzadko dysponuje pełną informacją o tym, jakie cechy profilu mu przypisano i jakie wagi mają poszczególne sygnały. Prawo do niepodlegania zautomatyzowanemu profilowaniu (art.~22 RODO \cite{rodo}) jest trudne do wyegzekwowania, gdy profilowanie jest wbudowane w podstawowy mechanizm działania usługi, a nie jedynie opcjonalną, wyraźnie oddzieloną funkcją \cite{netflix_recs, ig_ranking, amazon_privacy}.

Konieczne jest zatem nie tylko formalne zapewnienie możliwości skorzystania z praw podmiotów danych, ale też \textbf{projektowanie systemów, które czynią te prawa realnie dostępnymi} \cite{edpb2019}: panel wyjaśnień algorytmicznych, łatwa modyfikacja i resetowanie profilu rekomendacyjnego \cite{ig_reset, netflix_stop}, proste i trwałe mechanizmy wyboru feedu nieopartego na profilowaniu \cite{dsa}.

\subsection{Rola regulacji: RODO, DSA i przyszłe wyzwania}

Badane przypadki ilustrują zarówno siłę, jak i ograniczenia obowiązujących ram regulacyjnych. RODO \cite{rodo} ustanawia ważne zasady i prawa, ale ich egzekwowanie wobec bardzo dużych platform jest procesem długotrwałym i kosztownym -- postępowanie wobec Amazona w Luksemburgu toczyło się przez pięć lat i zakończyło uchyleniem kary przez sąd. DSA \cite{dsa} stanowi uzupełnienie RODO w zakresie przejrzystości systemów rekomendacyjnych \cite{ec_dsa}, ale jego praktyczne stosowanie wobec platform takich jak Meta (Instagram) jest wciąż w fazie kształtowania się \cite{ec_meta_dsa, reuters_ie}.

Nowe wyzwania regulacyjne przyniosą: integracja funkcji generatywnej AI z systemami rekomendacyjnymi (przypadek Meta AI w Instagramie \cite{meta_ai}), wzrastające możliwości łączenia danych z różnych ekosystemów \cite{himeur2022}, rekomendacje oparte na danych biometrycznych i emocjonalnych, oraz rosnąca złożoność łańcuchów przetwarzania danych obejmujących wielu podwykonawców \cite{rodo}.

\subsection{Konkluzja: ku odpowiedzialnemu projektowaniu systemów AI}

Opierając się na przeprowadzonych analizach, można sformułować kilka zasad \textbf{odpowiedzialnego projektowania systemów rekomendacyjnych} \cite{edpb2019, chen2025}:

\begin{enumerate}
  \item \textbf{Zasada proporcjonalności danych}: zakres danych wejściowych systemu powinien być ograniczony do tego, co jest rzeczywiście konieczne do realizacji funkcji rekomendacji, a nie do tego, co jest dostępne i technicznie możliwe do przetworzenia \cite{rodo}.
  \item \textbf{Zasada separacji celów}: dane używane do rekomendacji treści powinny być wyraźnie oddzielone od danych używanych do reklamy, analityki i celów wtórnych, a użytkownik powinien dysponować oddzielnymi kontrolami dla każdego z tych obszarów \cite{rodo, edpb2019}.
  \item \textbf{Zasada wyjaśnialności}: użytkownik powinien mieć dostęp do zrozumiałego wyjaśnienia, dlaczego widzi konkretne rekomendacje, jakie cechy profilu mu przypisano i jak może je zmodyfikować \cite{dsa, edpb2019}.
  \item \textbf{Zasada domyślnej minimalizacji}: system powinien być skonfigurowany tak, aby w stanie domyślnym przetwarzać możliwie najmniej danych, a rozszerzenie zakresu personalizacji powinno wymagać świadomej i prostej decyzji użytkownika \cite{rodo, edpb2019}.
  \item \textbf{Zasada ograniczonej retencji}: szczegółowe dane behawioralne powinny być przechowywane tylko tak długo, jak jest to rzeczywiście konieczne dla działania usługi i bezpieczeństwa, po czym powinny być agregowane, pseudonimizowane lub usuwane \cite{rodo, himeur2022}.
  \item \textbf{Zasada autonomii użytkownika}: projektowanie interfejsu powinno aktywnie wspierać, a nie utrudniać, korzystanie z mniej spersonalizowanych trybów działania; wybory prywatności powinny być pamiętane i respektowane, a dark patterns eliminowane \cite{dsa, reuters_nl}.
\end{enumerate}

Wdrożenie tych zasad wymaga nie tylko dobrej woli platform, ale przede wszystkim egzekwowalnych regulacji prawnych \cite{rodo, dsa}, aktywnych organów nadzorczych, niezależnych audytów algorytmicznych \cite{ec_meta_dsa} i świadomości użytkowników co do mechanizmów, które codziennie kształtują ich cyfrowe doświadczenie \cite{chen2025}.

% =========================================================
%  BIBLIOGRAFIA
% =========================================================
\newpage
\section*{Bibliografia}
\addcontentsline{toc}{section}{Bibliografia}

\begin{thebibliography}{99}

\bibitem{chen2025}
Chen, T., Liu, F., Shen, X.-L., Wu, J., Liu, Y. (2025). Conceptualization of privacy concerns and their influence on consumers' resistance to AI-based recommender systems in e-commerce. \textit{Industrial Management \& Data Systems}, 125(5), 1844--1868.

\bibitem{himeur2022}
Himeur, Y., Sohail, S.S., Bensaali, F., Amira, A., Alazab, M. (2022). Latest trends of security and privacy in recommender systems: A comprehensive review and future perspectives. \textit{Computers \& Security}, 118, 102746.

\bibitem{narayanan2008}
Narayanan, A., Shmatikov, V. (2008). Robust De-anonymization of Large Sparse Datasets. \textit{IEEE Symposium on Security and Privacy 2008}. \url{https://www.cs.cornell.edu/~shmat/shmat_oak08netflix.pdf}

\bibitem{rodo}
Rozporządzenie Parlamentu Europejskiego i Rady (UE) 2016/679 z dnia 27 kwietnia 2016 r. (RODO/GDPR). \url{https://eur-lex.europa.eu/eli/reg/2016/679/oj}

\bibitem{dsa}
Rozporządzenie Parlamentu Europejskiego i Rady (UE) 2022/2065 z dnia 19 października 2022 r. (DSA -- Akt o Usługach Cyfrowych). \url{https://eur-lex.europa.eu/legal-content/PL/TXT/?uri=CELEX\%3A32022R2065}

\bibitem{edpb2019}
European Data Protection Board. \textit{Guidelines 4/2019 on Article 25 Data Protection by Design and by Default} (wersja 2.0, 2019). \url{https://www.edpb.europa.eu/our-work-tools/our-documents/guidelines/guidelines-42019-article-25-data-protection-design-and_en}

\bibitem{netflix_privacy}
Netflix. \textit{Privacy Statement}. \url{https://help.netflix.com/legal/privacy}

\bibitem{netflix_recs}
Netflix Help Center. \textit{How Netflix's Recommendations System Works}. \url{https://help.netflix.com/en/node/100639}

\bibitem{netflix_data}
Netflix Help Center. \textit{What personal information Netflix holds about you}. \url{https://help.netflix.com/en/node/100624}

\bibitem{netflix_stop}
Netflix Help Center. \textit{How to stop certain uses of your personal information}. \url{https://help.netflix.com/en/node/100637}

\bibitem{netflix_transfer}
Netflix Help Center. \textit{Profile transfers}. \url{https://help.netflix.com/en/node/122698}

\bibitem{netflix_research}
Netflix Research. \textit{Personalization, Recommendations and Search}. \url{https://research.netflix.com/research-area/recommendations}

\bibitem{ig_ranking}
Instagram. \textit{Instagram Ranking Explained | How Our Algorithm Works}. \url{https://about.instagram.com/blog/announcements/instagram-ranking-explained}

\bibitem{ig_creators}
Instagram Creators. \textit{Understand how Instagram algorithms work}. \url{https://creators.instagram.com/grow/algorithms-and-ranking}

\bibitem{ig_recs}
Instagram Help Center. \textit{Recommendations on Instagram}. \url{https://help.instagram.com/313829416281232}

\bibitem{ig_reset}
Instagram Help Center. \textit{Reset your suggested content}. \url{https://help.instagram.com/556617736965724/}

\bibitem{meta_privacy}
Meta. \textit{Privacy Policy}. \url{https://www.facebook.com/privacy/policy/}

\bibitem{meta_ai}
Meta. \textit{Improving Your Recommendations on Our Apps With AI at Meta} (2025). \url{https://about.fb.com/news/2025/10/improving-your-recommendations-apps-ai-meta/}

\bibitem{meta_teen}
Meta. \textit{Instagram Expands Teen Accounts} (2026). \url{https://about.fb.com/news/2026/04/instagram-expands-teen-accounts-inspired-by-13-content-ratings/}

\bibitem{ec_dsa}
European Commission. \textit{Digital Services Act -- user rights}. \url{https://digital-strategy.ec.europa.eu/en/factpages/user-rights-under-digital-services-act}

\bibitem{reuters_nl}
Reuters. \textit{Dutch court orders Meta to change Facebook and Instagram timeline settings} (2025). \url{https://www.reuters.com/technology/dutch-court-orders-meta-change-facebook-instagram-timeline-settings-2025-10-02/}

\bibitem{reuters_ie}
Reuters. \textit{Ireland probes Meta's Instagram, Facebook over EU manipulation concerns} (2026). \url{https://www.reuters.com/legal/litigation/ireland-probles-metas-instagram-facebook-over-eu-manipulation-concerns-2026-05-05/}

\bibitem{ec_meta_dsa}
European Commission. \textit{Commission preliminarily finds Meta in breach of Digital Services Act} (2026). \url{https://ec.europa.eu/commission/presscorner/detail/en/ip_26_920}

\bibitem{amazon_privacy}
Amazon. \textit{Amazon.com Privacy Notice}. \url{https://www.amazon.com/gp/help/customer/display.html?nodeId=GX7NJQ4ZB8MHFRNJ}

\bibitem{amazon_dsp}
Amazon Ads. \textit{Amazon DSP}. \url{https://advertising.amazon.com/solutions/products/amazon-dsp}

\end{thebibliography}

\end{document}
